% =====================================================================
%  Planeación didáctica — Aprendizaje No Supervisado (FMAT-MIA-ANS)
%  Basada en Carta_Descriptiva-Aprendizaje_No_Supervisado.md
%  Formato tomado del modelo PD.pdf
%  Compilar con: pdflatex (dos veces) o latexmk -pdf
% =====================================================================
\documentclass[11pt,letterpaper]{article}

% ---------- Idioma y codificación ----------
\usepackage[utf8]{inputenc}
\usepackage[T1]{fontenc}
\usepackage[spanish,es-noshorthands,es-tabla]{babel}

% ---------- Tipografía (sans serif, como el modelo) ----------
\usepackage{helvet}
\renewcommand{\familydefault}{\sfdefault}
\usepackage{microtype}

% ---------- Página ----------
\usepackage[letterpaper,margin=2cm,top=1.8cm]{geometry}
\setlength{\parindent}{0pt}
\setlength{\parskip}{6pt}
\pagestyle{plain}

% ---------- Colores y tablas ----------
\usepackage[table]{xcolor}
\usepackage{array,longtable,tabularx}
\usepackage{enumitem}
\usepackage{titlesec}
\usepackage{needspace}
\usepackage[hidelinks]{hyperref}

\definecolor{marino}{HTML}{1F3864}   % títulos y encabezados de tabla
\definecolor{azul}{HTML}{2E75B6}     % subtítulo y reglas
\definecolor{texto}{HTML}{333333}
\definecolor{reqfondo}{HTML}{EAF2FA}
\definecolor{resfondo}{HTML}{EEF6EE}
\definecolor{filaclara}{HTML}{F5F9FD}

% Colores por bloque del calendario (oscuro / claro, alternados)
\definecolor{introA}{HTML}{E7E6E6}
\definecolor{uniUnoA}{HTML}{D6E7F2} \definecolor{uniUnoB}{HTML}{EAF3F9}
\definecolor{uniDosA}{HTML}{FFE3B8} \definecolor{uniDosB}{HTML}{FFF2DC}
\definecolor{uniTresA}{HTML}{F8C6DE} \definecolor{uniTresB}{HTML}{FCE4EF}
\definecolor{proyA}{HTML}{D5EAD2} \definecolor{proyB}{HTML}{E9F5E7}

\color{texto}

% ---------- Encabezados de sección ----------
\titleformat{\section}
  {\color{marino}\large\bfseries}
  {}{0pt}{\MakeUppercase}[{\color{marino}\titlerule[0.6pt]}]
\titlespacing*{\section}{0pt}{18pt}{10pt}

% ---------- Utilidades ----------
\newcolumntype{L}[1]{>{\raggedright\arraybackslash}p{#1}}
\newcommand{\C}[1]{\centering #1\\}
\newcommand{\bul}{\textbullet~}
% Texto blanco en negritas para encabezados de tabla
\newcommand{\hd}[1]{\leavevmode\color{white}\textbf{#1}}
% Celda de unidad: \unidad{Unidad 1}{Reducción de dimensionalidad}
\newcommand{\unidad}[2]{\leavevmode{\color{marino}\textbf{#1}}\newline{\footnotesize\itshape #2}}
% Celda de actividad (cursiva, como el modelo)
\newcommand{\act}[1]{{\itshape #1}}

% =====================================================================
\begin{document}

% ---------- Portada / encabezado ----------
\begin{center}
  {\color{marino}\fontsize{28}{32}\selectfont\bfseries Aprendizaje No Supervisado}\\[10pt]
  {\color{azul}\Large\itshape Planeación didáctica}\\[10pt]
  {\large 14 sesiones \;$\cdot$\; 2 horas por sesión \;$\cdot$\; 28 horas presenciales}\\[-4pt]
  {\color{azul}\rule{\textwidth}{0.8pt}}
\end{center}

% =====================================================================
\section{Descripción del curso}

Este curso ofrece a los estudiantes una introducción práctica al aprendizaje automático no supervisado: el descubrimiento de estructura en datos que no tienen etiquetas. Se cubrirán dos metodologías básicas --algoritmos de agrupamiento y reducción de la dimensionalidad-- así como el tema de detección de anomalías como una aplicación práctica. Cada sesión combina fundamentos conceptuales con implementación práctica en Python (scikit-learn, pandas, NumPy, Matplotlib).

\textbf{Competencia de la asignatura.} Implementa modelos de aprendizaje no
supervisado mediante técnicas estadísticas y algoritmos computacionales para la
resolución de problemas de sectores productivos y sociales, con evidencia de
rigurosidad técnica y sensibilidad respecto a su impacto en la toma de
decisiones informadas.

% =====================================================================
\section{Requisitos previos y resultados de aprendizaje}

\begin{center}
\renewcommand{\arraystretch}{1.25}
\begin{tabular}{|>{\columncolor{reqfondo}}p{0.44\textwidth}|>{\columncolor{resfondo}}p{0.48\textwidth}|}
\hline
{\color{marino}\textbf{Requisitos previos}}
\begin{itemize}[leftmargin=1.4em,topsep=2pt,itemsep=0pt]
  \item Programación introductoria en Python
  \item Estadística y probabilidad básicas
  \item Fundamentos de álgebra lineal
  \item Familiaridad con NumPy / pandas (recomendable)
\end{itemize}
&
{\color{marino}\textbf{Resultados de aprendizaje}}
\begin{itemize}[leftmargin=1.4em,topsep=2pt,itemsep=0pt]
  \item Preparar, escalar y explorar conjuntos de datos reales sin etiquetas
  \item Construir, comparar y validar modelos de agrupamiento
  \item Aplicar e interpretar técnicas de reducción de dimensionalidad
  \item Detectar anomalías con enfoques estadísticos, de densidad y basados en modelos
  \item Entregar un flujo de aprendizaje no supervisado completo (de extremo a
        extremo) y comunicar su impacto en la toma de decisiones
\end{itemize}
\\ \hline
\end{tabular}
\end{center}

{\footnotesize\itshape La carta descriptiva no establece requisitos académicos
previos; los anteriores son recomendaciones para seguir el ritmo del curso.}

% =====================================================================
\needspace{8\baselineskip}
\section{Evaluación}

\begin{center}
\renewcommand{\arraystretch}{1.5}
\rowcolors{2}{filaclara}{white}
\begin{tabularx}{\textwidth}{|X|c|}
\hline
\rowcolor{marino}
{\hd{Elemento de evaluación}} & {\hd{Peso}} \\ \hline
Miniproyecto -- Agrupamiento (sesión 6)                               & \textbf{30 \%} \\ \hline
Miniproyecto -- Reducción de dimensionalidad (sesión 10)                & \textbf{35 \%} \\ \hline
Miniproyecto -- Detección de anomalías (sesión 14)                     & \textbf{35 \%} \\ \hline
\end{tabularx}
\end{center}


% =====================================================================
\vspace{4pt}
\section{Calendario de sesiones}

{\small
\renewcommand{\arraystretch}{1.25}
\setlength{\tabcolsep}{4pt}
\begin{longtable}{|p{0.75 cm}|L{2.3cm}|L{3.0cm}|L{6.2cm}|L{3.0cm}|}
\hline
\rowcolor{marino}
{\centering \hd{\#}\\} &
{\hd{Unidad}} &
{\hd{Título de la sesión}} &
\centering{\hd{Temas principales}} &
{\hd{Laboratorio / Actividad}} \\ \hline
\endfirsthead
\hline
\rowcolor{marino}
{\centering \hd{\#}\\} &
{\hd{Unidad}} &
{\hd{Título de la sesión}} &
\centering{\hd{Temas principales}} &
{\hd{Laboratorio / Actividad}} \\ \hline
\endhead

% ---------------- Introducción ----------------
\rowcolor{introA}
\C{1} & \unidad{Introducción}{Fundamentos} &
\textbf{Introducción al aprendizaje no supervisado} &
\bul Aprendizaje supervisado vs.\ no supervisado; tipos de problemas\newline
\bul Preprocesamiento: valores faltantes, escalamiento y estandarización\newline
\bul Medidas de distancia y similitud (euclidiana, Manhattan, coseno, Mahalanobis)\newline
\bul La maldición de la dimensionalidad &
\act{Configuración del entorno; exploración de un conjunto de datos sin etiquetas} \\ \hline

% ---------------- Unidad 1 ----------------
\rowcolor{uniUnoA}
\C{2} & \unidad{Unidad 1}{Algoritmos de agrupamiento} &
\textbf{k-medias (\emph{k-means})} &
\bul Algoritmo de Lloyd y función objetivo (inercia)\newline
\bul Inicialización k-means++ y múltiples reinicios\newline
\bul Elección de $k$: método del codo y coeficiente de silueta\newline
\bul Limitaciones: forma esférica, escala y valores atípicos &
\act{Segmentación de clientes con k-medias} \\ \hline

\rowcolor{uniUnoB}
\C{3} & \unidad{Unidad 1}{Algoritmos de agrupamiento} &
\textbf{Agrupamiento jerárquico} &
\bul Enfoques aglomerativo y divisivo\newline
\bul Criterios de enlace: simple, completo, promedio y Ward\newline
\bul Dendrogramas y corte del árbol\newline
\bul Complejidad computacional y escalabilidad &
\act{Agrupamiento jerárquico con SciPy e interpretación de dendrogramas} \\ \hline

\rowcolor{uniUnoA}
\C{4} & \unidad{Unidad 1}{Algoritmos de agrupamiento} &
\textbf{Modelos de mezclas gaussianas (GMM)} &
\bul Agrupamiento probabilístico y asignaciones suaves\newline
\bul Algoritmo de esperanza-maximización (EM)\newline
\bul Tipos de matriz de covarianza\newline
\bul Selección de modelo con BIC y AIC &
\act{Comparación de GMM y k-medias en datos con grupos elípticos} \\ \hline

\rowcolor{uniUnoB}
\C{5} & \unidad{Unidad 1}{Algoritmos de agrupamiento} &
\textbf{Agrupamiento por densidad y espectral} &
\bul DBSCAN: puntos núcleo, frontera y ruido; parámetros $\varepsilon$ y \texttt{min\_samples}\newline
\bul HDBSCAN y OPTICS\newline
\bul Agrupamiento espectral: grafo de similitud y laplaciano\newline
\bul Grupos no convexos y de densidad variable &
\act{Comparación de DBSCAN, HDBSCAN y agrupamiento espectral en formas no convexas} \\ \hline

\rowcolor{uniUnoA}
\C{6} & \unidad{Unidad 1}{Algoritmos de agrupamiento} &
\textbf{Validación e interpretación de agrupamientos} &
\bul Índices internos: silueta, Davies-Bouldin, Calinski-Harabasz\newline
\bul Índices externos: ARI, información mutua normalizada\newline
\bul Estabilidad de los agrupamientos (remuestreo)\newline
\bul Perfilado de grupos y comunicación de resultados &
\act{Entrega del miniproyecto 1} \\ \hline

% ---------------- Unidad 2 ----------------
\rowcolor{uniDosA}
\C{7} & \unidad{Unidad 2}{Reducción de dimensionalidad} &
\textbf{Análisis de componentes principales (PCA)} &
\bul Varianza, covarianza y direcciones de máxima varianza\newline
\bul Eigendescomposición de la matriz de covarianza\newline
\bul Descomposición en valores singulares (SVD)\newline
\bul PCA como minimización del error de reconstrucción &
\act{Implementación de PCA desde cero con NumPy} \\ \hline

\rowcolor{uniDosB}
\C{8} & \unidad{Unidad 2}{Reducción de dimensionalidad} &
\textbf{PCA en la práctica y PCA con kernel} &
\bul Varianza explicada, gráfico de sedimentación y elección de componentes\newline
\bul Interpretación de cargas y biplots\newline
\bul PCA con kernel para estructuras no lineales\newline
\bul PCA incremental y aleatorizado para datos grandes &
\act{PCA con scikit-learn sobre imágenes de dígitos y Fashion-MNIST} \\ \hline

\rowcolor{uniDosA}
\C{9} & \unidad{Unidad 2}{Reducción de dimensionalidad} &
\textbf{Aprendizaje de variedades (\emph{manifold learning})} &
\bul Escalamiento multidimensional (MDS) e Isomap\newline
\bul t-SNE: similitudes probabilísticas y perplejidad\newline
\bul Riesgos al interpretar proyecciones no lineales &
\act{Comparación visual de PCA, MDS, Isomap y t-SNE sobre Fashion-MNIST} \\ \hline

\rowcolor{uniDosB}
\C{10} & \unidad{Unidad 2}{Reducción de dimensionalidad} &
\textbf{Modelos de variables latentes y autoencoders} &
\bul Análisis factorial y PCA probabilístico\newline
\bul Análisis de componentes independientes (ICA)\newline
%\bul Autoencoders: codificador, decodificador y espacio latente\newline
\bul Comparación de métodos lineales y no lineales &
\act{Entrega del miniproyecto 2} \\ \hline


% ---------------- Unidad 3 ----------------
\rowcolor{uniTresA}
\C{11} & \unidad{Unidad 3}{Detección de anomalías} &
\textbf{Enfoques estadísticos y basados en densidad} &
\bul Tipos de anomalías: puntuales, contextuales y colectivas\newline
\bul Puntuación $z$, IQR y distancia de Mahalanobis\newline
\bul Estimación de densidad con kernels y GMM\newline
\bul Factor de valor atípico local (LOF) &
\act{planteamiento del miniproyecto 3: Detección de anomalías en conjuntos de ODDS.} \\ \hline

\rowcolor{uniTresB}
\C{12} & \unidad{Unidad 3}{Detección de anomalías} &
\textbf{Métodos basados en modelos y su evaluación (I)} &
\bul Bosque de aislamiento (\emph{Isolation Forest}) \newline
\bul SVM de una clase (\emph{One-Class SVM}) &
\act{Comparativa de detectores con PyOD en un conjunto de fraude con tarjetas} \\ \hline

\rowcolor{uniTresA}
\C{13} & \unidad{Unidad 3}{Detección de anomalías} &
\textbf{Métodos basados en modelos y su evaluación (II)} &
\bul Error de reconstrucción con PCA y autoencoders \newline
\bul Evaluación con pocas etiquetas: precisión@k, curva PR, ROC-AUC &
\act{Comparativa de detectores con PyOD en un conjunto de fraude con tarjetas} \\ \hline

% -------------- Flujo integrador y Cierre ------------------
\rowcolor{uniTresB}
\C{14} & \unidad{Unidad 3}{Detección de anomalías} &
\textbf{Flujo integrador y cierre del curso} &
\bul Diseño del flujo completo: preparación, reducción, agrupamiento y anomalías\newline
\bul Selección e interpretación de métricas\newline
\bul Impacto de los resultados en la toma de decisiones\newline
\bul Ética y sesgos en el análisis no supervisado \newline
\bul Síntesis del curso y temas avanzados &
\act{Entrega del miniproyecto 3} \\ \hline

\end{longtable}
}

% =====================================================================
\section{Software y recursos}

\begin{center}
\renewcommand{\arraystretch}{1.2}
\rowcolors{1}{filaclara}{white}
\begin{tabularx}{\textwidth}{|>{\raggedright\arraybackslash\leavevmode\color{marino}\bfseries}p{4.2cm}|X|}
\hline
Lenguaje de programación & Python 3.10+ \\ \hline
Bibliotecas principales  & scikit-learn, pandas, NumPy, SciPy, Matplotlib, seaborn \\ \hline
Bibliotecas avanzadas    & PyOD, PyTorch (autoencoders) \\ \hline
Entorno de desarrollo    & Jupyter Notebook / JupyterLab, VS Code con la extensión de Python o Google Colab \\ \hline
Bibliografía básica      &
James, G., Witten, D., Hastie, T., Tibshirani, R. y Taylor, J. (2023).
\emph{An Introduction to Statistical Learning: With Applications in Python}.
Springer. ISBN 978-3-031-38747-0. (Cap.~12)\par\vspace{3pt}
Bishop, C. M. (2006). \emph{Pattern Recognition and Machine Learning}.
Springer. ISBN 978-0387310732. (Caps.~9 y 12)\par\vspace{3pt}
Murphy, K. P. (2022). \emph{Probabilistic Machine Learning: An Introduction}.
MIT Press. ISBN 978-0262046824. (Caps.~20 y 21) \\ \hline
Bibliografía complementaria &
Hinton, G. E. y Sejnowski, T. J. (Eds.) (1999). \emph{Unsupervised Learning:
Foundations of Neural Computation}. MIT Press. ISBN 978-0262581684. \\ \hline
Documentación en línea   &
Guía de usuario de scikit-learn (\emph{Unsupervised learning}); documentación de PyOD \\ \hline
Conjuntos de datos       &
\texttt{sklearn.datasets}, Fashion-MNIST, Repositorio UCI ML, ODDS (Outlier
Detection DataSets), Kaggle, datos abiertos de INEGI y datos.gob.mx \\ \hline
\end{tabularx}
\end{center}

{\footnotesize\itshape\color{gray}
Nota: Los temas de las sesiones pueden ajustarse al ritmo del grupo. Todos los
archivos de laboratorio y los conjuntos de datos se distribuirán a través de la
plataforma del curso.}\\[-6pt]
{\color{gray}\rule{\textwidth}{0.4pt}}

\end{document}
